\documentclass[11pt,letterpaper]{article}
\usepackage[margin=1in]{geometry}
\usepackage{amsmath}
\usepackage{mathptmx}
\usepackage[T1]{fontenc}
\usepackage{enumitem}
\usepackage{titlesec}
\setlist{nosep,leftmargin=1.5em}
\titlespacing*{\section}{0pt}{8pt}{3pt}
\titleformat{\section}{\normalfont\bfseries}{\thesection}{0.5em}{}
\setlength{\parskip}{4pt}
\setlength{\parindent}{0pt}
\pagestyle{plain}

\begin{document}

\begin{center}
{\large\bfseries CPTS 360 Programming Assignment 2: CPU Scheduling Report}\\[2pt]
Isaiah Foster \quad WSU ID: XXXXXXXX
\end{center}

\section{Overall Design}

The simulator (\texttt{scheduler.c}) reads \(n\) and the \((A,B,C,M)\) quadruples from the file named on the
command line. Each quadruple is parsed with \texttt{fscanf(" ( \%u \%u \%u \%u )")}, so any amount of
whitespace is accepted inside and between the parentheses, and anything after the \(n\)th quadruple (the
``comments'') is never read. Every process gets a \texttt{processID} equal to its position in the input.

All three algorithms share one function, \texttt{simulate(algorithm, \ldots)}, which runs on a fresh copy of the
input. It differs between algorithms in only two places: whether the running process can be preempted (RR only)
and how the next process is taken from the ready queue (FIFO for FCFS and RR, minimum remaining time for SJF).
Each process stores a state (\textit{unstarted}, \textit{ready}, \textit{running}, \textit{blocked},
\textit{terminated}), the remaining CPU burst, the remaining I/O burst, the length of its last CPU burst, the RR
quantum left, and counters for CPU time, I/O time and waiting time.

Time advances in integer cycles; cycle \(t\) covers \([t, t+1)\). Each iteration of the main loop does the
following:
\begin{enumerate}
  \item \textbf{Running process transitions.} If the process on the CPU has run for \(C\) cycles in total, it
        terminates with finishing time \(t\) and is appended to the finished list. Otherwise, if its CPU burst is
        exhausted, it blocks for \(\text{burst}\times M\) cycles. Otherwise, under RR, if its quantum is exhausted
        it is preempted back to ready.
  \item \textbf{Blocked and arriving processes.} Blocked processes whose I/O burst has reached zero, and
        unstarted processes with \(A = t\), become ready.
  \item \textbf{Ready queue insertion.} Every process that became ready in this cycle, whether it was preempted,
        finished I/O or just arrived, is gathered into a batch. The batch is sorted by the tie-breaking rule and
        then appended to the ready queue.
  \item \textbf{Dispatch.} If the CPU is idle and the ready queue is non-empty, a process is selected. If it has no
        CPU burst left over from a preemption, a new burst is drawn with \texttt{randomOS(B)} and capped at the
        CPU time remaining (\(C\) minus cycles run). The RR quantum is reset to 2.
  \item \textbf{Accounting.} The running process gains one CPU cycle, and its burst and quantum each drop by
        one. Each blocked process gains one I/O cycle, and each ready process gains one waiting cycle. If any
        process is blocked, the global count of I/O-busy cycles goes up by one. Then \(t\) is incremented.
\end{enumerate}
The loop ends when all \(n\) processes have terminated. The last termination happens at the start of a cycle,
and the counter is still incremented after it, so the overall finishing time is \texttt{CURRENT\_CYCLE - 1}, as
the provided \texttt{printSummaryData()} expects. CPU utilization is total CPU time divided by the finishing
time. I/O utilization is the number of cycles in which at least one process was blocked, divided by the
finishing time.

\section{The Three Algorithms}

\textbf{FCFS.} The ready queue is a FIFO array of process indices. Dispatch always takes the head, and the running
process keeps the CPU until its burst ends (it blocks) or it terminates. A process that comes back from I/O
goes to the tail of the queue, so the CPU is shared in order of when each process became ready.

\textbf{RR (quantum 2).} This uses the same FIFO queue and dispatch as FCFS. On dispatch the quantum is set to 2,
and step 1 preempts the process when the quantum reaches zero and its burst is not yet finished. The preempted
process keeps its remaining CPU burst, so no new random number is drawn when it next runs. Its eventual I/O
burst is still computed from the full burst it originally drew (stored in \texttt{lastCPUBurst}).

\textbf{SJF.} This is not preemptive, but a process gives up the CPU when it blocks for I/O. Priority is the total
CPU time remaining, \(C - \text{cycles run}\), not the length of the next burst. On dispatch,
\texttt{dequeueReady()} scans the whole ready queue for the process with the smallest remaining time and
removes it. A running process is never interrupted by a shorter job that arrives later; the shorter job waits
until the CPU becomes free.

\section{Tie-Breaking}

\textbf{Simultaneous events for the running process.} Step 1 uses an \texttt{if / else if / else if} chain in
the order the spec requires: termination first, then blocking, then preemption. If a process's last CPU cycle
also ends its burst and its quantum, it terminates and does not block. If its burst and quantum end together,
it blocks and is not preempted.

\textbf{Processes with equal priority.} One comparison function, \texttt{favoredOver(a, b)}, implements the
rule: the earlier arrival time \(A\) wins, and if the arrival times are equal, the lower \texttt{processID} (the
process listed earlier in the input) wins. It is used in two places:
\begin{itemize}
  \item \textbf{FCFS and RR:} all processes that become ready in the same cycle (preempted, returning from I/O or
        newly arrived) are collected into one batch. The batch is sorted with \texttt{favoredOver} before it is
        appended to the FIFO queue, so simultaneous insertions are ordered by \((A, \text{input order})\).
        Processes already in the queue keep their earlier positions.
  \item \textbf{SJF:} when two ready processes have the same remaining CPU time, \texttt{favoredOver} decides
        between them during the minimum scan. This makes the choice independent of the order in which
        processes entered the queue.
\end{itemize}
The rule breaks every tie, because no two processes share both the same \(A\) and the same input position.
Only one process runs at a time, so at most one process can terminate in a given cycle, and the order of the
finished list is therefore also fully determined.

\section{Testing}

The program reproduces \texttt{sample\_io/output/summary/output-1} through \texttt{output-3} byte for byte. The
samples only use \(B = 1\), so I also checked RR preemption by hand on \texttt{2 (0 5 6 1) (0 5 6 1)}. That run
draws bursts longer than the quantum, has a burst and a quantum end in the same cycle, and has a preempted
process tie with a process returning from I/O; every metric matched the hand trace. Further inputs with
unsorted arrival times, irregular spacing, trailing comments and \(M = 0\) ran cleanly under AddressSanitizer
and UndefinedBehaviorSanitizer.

\end{document}
